\documentclass[a4paper,11pt,oneside]{article}
\usepackage{a4wide}                     % Iets meer tekst op een bladzijde
\usepackage[english,dutch]{babel}       % Voor nederlandstalige hyphenatie (woordsplitsing)
\usepackage{amsmath}                    % Uitgebreide wiskundige mogelijkheden
\usepackage{amssymb}                    % Voor speciale symbolen zoals de verzameling Z, R...
\usepackage{url}                        % Om url's te verwerken
\usepackage{graphicx}                   % Om figuren te kunnen verwerken
\usepackage[small,bf,hang]{caption}    % Om de captions wat te verbeteren
\usepackage{xspace}                     % Magische spaties na een commando
\usepackage[latin1]{inputenc}           % Om niet ascii karakters rechtstreeks te kunnen typen
\usepackage{float}                      % Om nieuwe float environments aan te maken. Ook optie H!
\usepackage{flafter}                    % Opdat floats niet zouden voorsteken
\usepackage{listings}                   % Voor het weergeven van letterlijke text en codelistings
\usepackage{marvosym}                   % Om het euro symbool te krijgen
\usepackage{textcomp}                   % Voor onder andere graden celsius
\usepackage[left=2.5cm,top=1.4cm,right=1.5cm,bottom=25mm,nohead]{geometry}   % marges instellen
\usepackage{helvet}		% We gebruiken Helvetica.  Arial is de Microsoft-versie van Helvetica.  Dit wijzigt de schreefloze letters in helvetica, de rest blijft zoals standaard.
\usepackage{fancyhdr}                   % Voor fancy headers en footers.
\pagestyle{fancy}                       % De bladzijdestijl
\usepackage{graphics}                   % Om figuren te verwerken.
\usepackage{lastpage}										% Om aan de laatste pagina te kunnen.
\newcommand{\npar}{\vspace{0.5cm} \par \noindent}      % witruimte aan paragrafen


% begin van het document
\begin{document}
\sffamily               % We gebruiken overal een schreefloze letter (hier dus Helvetica)


% Header (titel en datum aanpassen)
\fancyhf{}                              % Resetten fancy settings.
\renewcommand{\headrulewidth}{0pt}      % Geen lijn onder de header.
\hspace{-7mm}\begin{tabular}{p{50mm}c}
\includegraphics[width=40mm]{logo}& 
	\begin{tabular*}{120mm}[b]{c}
	\huge{Studentenvisie op het Gebruik} \\
	\huge{en de Kwaliteit van Slides} \\
	\vspace{1cm}\\
	\Large{30 februari 2005} \\
	\end{tabular*}
\end{tabular}
\vspace{1cm}


% De footer
{ 
\lfoot{ \sffamily 
        \vspace{1ex} \small
				Studentenvisie op het Gebruik en de Kwaliteit van Slides \\
        De Stuvers - Facultaire Raad van Studenten \\
        Keizer Karelstraat 38, 9000 Gent
}
\rfoot{ \sffamily
				\vspace{1ex} \small
				Pagina \thepage/\pageref{LastPage}\\
        \vspace{1.9ex} \small
        www.GSR.UGent.be
}
}\renewcommand{\footrulewidth}{0.1pt}


% De inhoud
\section{\sffamily Algemene uitgangspunten}
\npar Elektronische presentaties hebben de laatste jaren enorm veel aan belang gewonnen.  Meer en meer lesgevers schakelen van het klassieke bord werk over naar een kant-en-klare presentatie.  Door dit op voorhand voor te bereiden, heeft men tijdens de hoorcolleges de tijd om op een betere manier de uitleg te verzorgen zonder constant (bord)schrijf werk te moeten verrichten.  Goed uitgewerkte figuren in de presentatie bieden ook een niet te onderschatten meerwaarde in het verduidelijken van de leerstof. Worden deze slides ook aan de studenten aangeboden (hetzij via Minerva, hetzij in syllabus vorm), dan kunnen de studenten tijdens de presentatie, de les beter mee volgen in plaats van constant het bord te moeten overpennen.  Verder zijn slides ook uitstekend te gebruiken voor sterk vernieuwende (of nieuwe) cursussen, zoals bijvoorbeeld vaak het geval is in de opleiding computerwetenschappen waar het al eens over state-of-the-art technieken of in ontwikkeling zijnde technologie�n gaat.  Wij staan dus helemaal niet weigerachtig tegenover het gebruik van elektronische presentaties in de lessen, integendeel, ze vormen zeker een toegevoegde waarde.
\npar Het kan echter wel een probleem worden wanneer de slides deel uitmaken van de leerstof en de cursus aanvullen of, in extremis, een cursus volledig vervangen en zelf voor syllabus doorgaan.  Het is dan uitermate belangrijk dat de presentaties voor iedereen vlot begrijpbaar en goed leesbaar zijn.  Weinigen studeren graag achter een computerscherm, waardoor de kwaliteit van de slides ook na afdrukken moet behouden blijven.  Zeker bij vakken waar er enkel slides zijn als cursusmateriaal en deze gebruikt mogen worden tijdens een open-boek examen, is de goede afdrukbaarheid van het allergrootste belang en zelfs een noodzaak.  Hieronder sommen we een aantal tips op om slides zo goed mogelijk te laten voldoen aan de leesbaarheids- en afdrukbaarheidsvereisten.

\section{\sffamily Voorstellen en raadgevingen}
\npar Deze voorstellen zijn vooral van toepassing op de slides die de student kan afdrukken (bijvoorbeeld vanop Minerva) of afgedrukt krijgt. Worden de slides enkel als illustratie tijdens de lessen gebruikt, dan wordt het uiteraard aangemoedigd om deze raadgevingen te volgen maar zijn ze niet zo cruciaal.
\begin{itemize}
\item Nummer de slides.
\item Gebruik een donkere tekst op een lichte, sobere (of liefst: geen) achtergrond.
\item Zorg ervoor dat de slides ook in zwart-wit goed leesbaar zijn, want weinig studenten drukken hun cursussen af met een kleurenprinter. Als er bijvoorbeeld een kader rond een tekst of formule geplaatst wordt om het belang ervan aan te duiden, gebruik dan bij voorkeur geen achtergrondkleur in dat kader.
\item Gebruik geen exotische lettertypes.  Studenten die het bewuste lettertype niet op hun PC staan hebben, kunnen de slides niet (goed) lezen.  Dit kan ook vermeden worden door ook een .pdf-versie van de slides ter beschikking te stellen indien deze goed aangemaakt is (zie verder).
\item Gebruik geen semi-transparantie (zoals in PowerPoint XP). Bij studenten die een oudere versie van PowerPoint gebruiken, verdwijnt deze transparantie en wordt het object volledig zwart (of gekleurd) en de bovenstaande tekst bijgevolg onleesbaar bij het afdrukken.
\item Plaats geen objecten boven elkaar.  Een tekst waarboven tekstballonnen of andere objecten verschijnen is geen goed idee.  De onderliggende informatie verdwijnt immers bij het afdrukken.  Let bijgevolg ook op met animaties en effecten.
\item De laatste versie van MS Office PowerPoint en het Open Office equivalent zijn op de Athena server te vinden. Het systeem vertoont echter nog teveel gebreken (zoals het regelmatig uitvallen van de verbinding) om er vanuit te gaan dat iedereen die versie kan gebruiken. Niet iedereen beschikt trouwens over een internetverbinding (bijvoorbeeld op kot).  Is de presentatie gemaakt in de laatste versie van PowerPoint, dan kunnen incompatibiliteiten optreden met bvb. de PowerPoint 2000 versie. De slides worden bijgevolg onleesbaar.  Het is dan ook goed om naast (of ter vervanging van) de .ppt-versie ook een .pdf-versie aan te maken en deze ter beschikking te stellen.  Dit kan op een uiterst eenvoudige manier door middel van bijvoorbeeld een virtuele printer waarnaar afgedrukt kan worden zoals PDFCreator (\url{http://pdfcreator.sf.net}), Adobe Acrobat Destiller of via een exporteerfunctie zoals in OpenOffice.org.  Op die manier krijgt iedereen de slides zoals deze bedoeld waren. Door de slides naar .pdf om te zetten kan ook meteen nagegaan worden in hoeverre de slides leesbaar zijn na afdrukken. 
\item Indien u een .pdf-versie ter beschikking stelt en deze aanmaakt door af te drukken naar een .pdf-printer, kies dan voor ��n slide per pagina, zo groot mogelijk.  Op die manier kan de student zelf kiezen hoeveel pagina's er op ��n blad geprint moeten worden en kan de fullscreen-functie van Acrobat Reader gebruikt worden om de presentatie te bekijken zoals ze in de les gegeven werd (zonder animaties uiteraard).
\item Slides worden best ook zo vroeg mogelijk (via Minerva) verspreid zodat notities tijdens de les direct bij de juiste slides kunnen geplaatst worden. Lukt het niet om ze tijdig online te plaatsen, dan is nummering van de slides onontbeerlijk!
\item Als de slides als cursus dienen, is de structuur en een accurate inhoudstafel zoals bij een klassieke syllabus noodzakelijk, zeker 
als niet alle slides in ��n keer uitgedeeld worden. 

\end{itemize}
\npar Tot slot zouden we er nog eens willen de nadruk op leggen dat de leesbaarheid van slides van het allergrootste belang is voor de studeerbaarheid van een vak, zeker als deze slides het enige lesmateriaal zijn.



\end{document}
